
Across 4,493 open-weight language models, naive OLS regression shows that models documenting their pretraining data curation practices score lower on standardized benchmarks. This result is traced to a size-vintage confound: well-documented models are predominantly older, smaller-parameter base models, while high-scoring V2 models tend to be newer, larger, or fine-tuned derivatives. The confound itself is a finding that motivates propensity-matched analysis as the appropriate next step.

The contributions of this work are:

\begin{enumerate}
\item \textbf{The LLM Documentation-Benchmark Registry (n = 4,493):} The first large-scale dataset linking binary curation documentation indicators to V2 benchmark scores, revealing that 96.5\% of evaluated models document zero curation practices.
\end{enumerate}

\begin{enumerate}
\item \textbf{Targeted family sampling:} A retrieval methodology ensuring sufficient feature variance (3 of 4 binary features) within API rate constraints, with explicit acknowledgment of the resulting selection bias.
\end{enumerate}

\begin{enumerate}
\item \textbf{Confound identification:} Documentation of the size-vintage confound behind beta = -3.45 and the perplexity filtering vocabulary gap.
\end{enumerate}

Several directions for future work are motivated by these findings. Bridge validation should test whether documented families show measurably different corpus hygiene metrics using models with accessible pretraining corpora. Propensity-matched regression should attempt to isolate documentation effects within parameter strata, controlling for model size and vintage. Benchmark-type differential analysis should reformulate the differential sensitivity claim for V2 benchmarks. Vocabulary-aware NLP extraction should expand feature coverage beyond canonical terminology.

Whether model cards contain genuine signals about training data quality remains an open question. The registry provides the infrastructure to make it an answerable one.
